\documentclass[preview]{standalone}

\usepackage{amsmath}
\usepackage{amssymb}
\usepackage{stellar}
\usepackage{definitions}

\begin{document}

\id{minimal-polynomial}
\genpage

\section{Polynomial evaluation}

\includesnpt{linearalgebra-polynomial-of-a-matrix-definition}

\begin{snippetdefinition}{matrix-minimal-polynomial-definition}{Minimal polynomial}
    Let \(A\in\matrices_n(K)\), where \(K\) is a \field. Its
    \emph{annihilating ideal} is
    \[
        \mathcal I_A
        =
        \{f\in K[X]\suchthat f(A)=0\}.
    \]
    The \emph{minimal polynomial} of \(A\), denoted by \(\mu_A\), is the
    unique monic generator of \(\mathcal I_A\).

    Equivalently, if \(A\) represents a \linearendomorphism
    \(\varphi\colon V\fromto V\), then \(\mu_\varphi=\mu_A\) is the
    monic polynomial of least degree such that
    \[
        \mu_\varphi(\varphi)=0,
        \qquad\text{that is,}\qquad
        \mu_\varphi(\varphi)(v)=0_V
        \quad\text{for every }v\in V.
    \]
\end{snippetdefinition}

\begin{snippetproposition}{matrix-annihilating-polynomials-ideal}{Annihilating polynomials form an ideal}
    For every \(A\in\matrices_n(K)\), the \set
    \[
        \mathcal I_A
        =
        \{f\in K[X]\suchthat f(A)=0\}
    \]
    is a nonzero \ideal of \(K[X]\). Consequently, it has a unique monic
    generator \(\mu_A\).
\end{snippetproposition}

\begin{snippetproof}{matrix-annihilating-polynomials-ideal-proof}{matrix-annihilating-polynomials-ideal}{Annihilating polynomials form an ideal}
    The zero polynomial belongs to \(\mathcal I_A\). If
    \(f,g\in\mathcal I_A\), then
    \[
        (f+g)(A)=f(A)+g(A)=0.
    \]
    If \(h\in K[X]\), then
    \[
        (hf)(A)=h(A)f(A)=0.
    \]
    Thus \(\mathcal I_A\idealin K[X]\).

    Regard \(K^n\) as the \(K[X]\)-module associated with
    \(v\mapsto Av\). This module is torsion. If
    \(e_1,\ldots,e_n\) is the standard basis, choose nonzero
    \(f_i\in K[X]\) with \(f_i(A)e_i=0\). Then the nonzero product
    \(f_1\cdots f_n\) annihilates every \(e_i\), hence belongs to
    \(\mathcal I_A\). Since \(K[X]\) is a
    \principalidealdomain, the nonzero ideal \(\mathcal I_A\) has a
    generator, unique after requiring it to be monic.
\end{snippetproof}

\begin{snippetcorollary}{minimal-polynomial-divisibility-property}{Divisibility property of the minimal polynomial}
    A polynomial \(f\in K[X]\) is an
    \annihilatingpolynomial for \(A\)
    \ifandonlyif
    \[
        \mu_A\divides f.
    \]
\end{snippetcorollary}

\begin{snippetproof}{minimal-polynomial-divisibility-property-proof}{minimal-polynomial-divisibility-property}{Divisibility property of the minimal polynomial}
    By definition,
    \[
        \mathcal I_A=(\mu_A).
    \]
    Thus \(f(A)=0\) exactly when \(f\in(\mu_A)\), which is equivalent to
    \(\mu_A\divides f\).
\end{snippetproof}

\section{Similarity invariance}

\begin{snippetproposition}{minimal-polynomial-similarity-invariance}{The minimal polynomial is invariant under similarity}
    If \(A,B\in\matrices_n(K)\) are similar, then
    \[
        \mu_A=\mu_B.
    \]
    Consequently, the minimal polynomial of a
    \linearendomorphism is independent of the basis used to represent
    it.
\end{snippetproposition}

\begin{snippetproof}{minimal-polynomial-similarity-invariance-proof}{minimal-polynomial-similarity-invariance}{The minimal polynomial is invariant under similarity}
    Write \(B=P^{-1}AP\). For every \(m\geq0\),
    \[
        B^m=P^{-1}A^mP.
    \]
    Therefore, for every \(f\in K[X]\),
    \[
        f(B)=P^{-1}f(A)P.
    \]
    Hence \(f(A)=0\) \ifandonlyif \(f(B)=0\), so
    \(\mathcal I_A=\mathcal I_B\). Their monic generators coincide.
\end{snippetproof}

\end{document}
